\documentclass[a4paper,11pt]{article}
\usepackage{jinstpub} % for details on the use of the package, please see the JINST-author-manual
\usepackage{lineno}
\usepackage{color}
\linenumbers

% Proceedings/Special Issues
% Please note that this macro will be edited in production 
%% \proceeding{N$^{\text{th}}$ Workshop on X\\
%% When\\
%% Where}



\title{\boldmath Ionization charge unfolding at zero-suppressed liquid argon time projection chamber}

% Collaborations

%% [A] If main author
%% \collaboration{\includegraphics[height=17mm]{collabroation-logo}\\[6pt]
%%  XXX collaboration}

%% or
%% [B] If "on behalf of"
%% \collaboration[c]{on behalf of XXX collaboration}


% Authors
% Please note that in JINST a corresponding author is required alongside with their e-mail addres
% The "\note" macro will give a warning: "Ignoring empty anchor...", you can safely ignore it.

%% [A] simple case: 2 authors, same institution
%% \author[1]{A. Uthor\note{Corresponding author.}}
%% \author{and A. Nother Author}
%% \affiliation{Institution,\\Address, Country}

%% or, e.g.
%% [B] more complex case: 4 authors, 3 institutions, 2 footnotes
%% \author[a,b,1]{F. Irst,\note{Corresponding author.}}
%% \author[c]{S. Econd,}
%% \author[a,2]{T. Hird\note{Also at Some University.}}
%% \author[c,2]{and Fourth}
%% \affiliation[a]{Institution_1,\\Address, Country}
%% \affiliation[b]{Institution_2,\\Address, Country}
%% \affiliation[c]{Institution_3,\\Address, Country}

\author{An Author}
\affiliation{Physics Department, Brookhaven National Laboratory,\\
Upton, NY, 11973 U.S.A}

% E-mail addresses: only for the corresponding author
\emailAdd{no.one@bnl.gov}

\abstract{The liquid argon time projection chamber at near site of the deep underground neutrino experimetn adopts a zero-suppressed pixelated readout detector. This expose challeges on traditional signal processing approach, which relies on the full equally-spaced sampling of the waveform. In the liquid argon TPC equipped with projective readout, ionization charges are unfolded by deconvoving the electronic readout response as well as the induced template by drifting charges. The zero suppression for small amplitudes destruct the foundation of traditional deconvolution approach. In this work, we present a new method to deal with charge unfolding by physics-guided waveform compensation. The charge regression performance is evaluated in terms of residues between true ionization charge and deconvolved charge. The residue width is shown to be below 500 electrons.
}

\keywords{Only keywords from JINST's keywords list please}

\arxivnumber{1234.56789} % Only if you have one

\begin{document}
\maketitle
\flushbottom

\section{Introduction}
\label{sec:intro}

The deep underground neutrino experiment (DUNE) is the next generation long-baseline neutrino experiment for measuring neutrino oscillation parameters including CP violation factor. The near detector at Fermilab consists of a modular set of liquid argon time projection (LArTPC), muon spectrameter, and System for on-Axis Neutrino Detection. The LArTPC at near site is referred to as ND-LAr. The near detector is designed to be . There are four modules at Sanford underground research facility. Multi-view projective readout is employed at far detector while pixelated design is adopted at near-site detector.

The NDLAr is only of $7\times 5 \times 3\times$ cubic meters large while the far detector module is of $60 \times 14 \times 14$ cubic meters. The total number of readout channels at FD is of $\mathcal{O}(5\times 10^5)$, while the corresponding number at ND-LAr is almost the same. The low density of channel allows the FD to use a sampling detector. The full waveform from a drifting charge cloud will be recorded for future signal processing. On the other hand, the high density at ND-LAr drives the zero-suppressing detector. Low charge activity are not recorded due to a threshold. A short integration window effectively serve as a low pass filter is used. The total recorded charge is assumed to be the amount of charge and trigger time is assumed to be time of arrival. This is a good approximation when the signal is sharp and slew is small among different signals.

This is not generally the case for a drifting charge when there is no screening before readout electronics. The signal from ionization charge is typically long and the worse case is it can generate currents on neighboring pixels to its collection point. The coherent sum from several sources can falsely trigger a channel. We can observe pre-triggering before the charge arrives the anode. It is necessary to eliminate the ghost hits from pretriggering so that ionization charge's amount and positions are regressed precisely. These are essential for neutrino energy measurements, vertex reconstruction, and flavor tagging, in the oscillation measurement.

Traditional signal processing requires the system is linear and time invariant. The detector response and readout electronics can then be written as a universal convolution kernel. In addition, the full history of the waveform must be preserved so that we can deconvolve the recorded waveform from the response function given the system. In the zero-suppressed system, the recorded signal is not linear anymore, the direct convolution on the raw measurements is invalid. We develop a new approach to extend the deconvolution technique to the ND-LAr in this article.

\section{Signal processing at NDLAr}
\subsection{Review of the traditional signal processing}
According to Shannon sampling theorem, the system information can be recovered when Nyquist frequency is over the half bandwidth the system. In nature, the system generally have a long tail in high frequency coupled with noises. A low pass filter is usually employed to eliminate the high frequency tail. Together with other electronic responses, the system can be written as $m=r*s+n$, where $r(t)$ consists of both contributions from the field response and sampling electronc responsem and $n(t)$ is the noise in the system. The decvonlution is the revser process, usually implemented through FFT. In the frequency domain, $\hat{S}(\omega)=M/R-N/R$. It generally common to adapt a Gaus filter to suppress the blow up at small $R(\omega)$ in the second noise term, while preserving the charge integral. On the other hand, the noise can be further eliminated by defining region of interests. Areas with small charges are considered as signal-free to be removed. Identifying the region of interests is done by enforcing Wiener-inspired filters. These filters naturally give us the resolution of charge in both space and amounts. In the reality, $R$ is usually position dependent. However, in the deconvolution, the ignore along the detector channel pitch limits us to use an average response from several paths. This brings additional smearing effects. In addition, because of the presence of the low-pass filter, the fine structure of the original ionization charge is blurred. Limited to all the factors, the reconstructed charge is the true ionization convolved with effective resolution factors.

\subsection{Review of readout logic at NDLAr}
The NDLAr adopts a self-triggering scheme. Electronics is dormnant for most of time. The charge is accumulated on its frontend, and discriminator monitors the charge at CSA.
As long as the accumulate charge exceeds the threshold, a trigger cycle is issued. After a configurable delayed time, the charge at frontend is recorded if it is still larger than the threshold value. This veto process is to reject random noises. It is worth noting that during the ADC digitization conversion time, trigger cycle is suppressed. Then the charge pre-trigger can only be known for the first triggering. We lose the exact time when charge passes triggering for consecutive hits. One time clock after the digization starts, the frontend resets and drain the charge. New charge accumulation starts. In this mechanism, negative contributions from neighboring charges are blind from the readout. And we may overestimate the charges from this process. On the other hand, the reset mechanism and leftover charge below threshold would be lost. Calibration of the charges must be performed topology-dependent.

There is also an variant operational mode for LArPix, namely, burst mode. It is an extension for the self-triggering mode. The frontend does not reset after passing the first delay time window, instead it keeps charge accumulation and continously record a few more charge amounts at a fixed time window. At the end of the burst sequence it resets the frontend again. The burst frequency is typicall chosen the same as the delay window length as in the self-trigger scheme, much slower than the system time clocking, which is used in trigger timestamp.

\subsection{Algorithm for signal processing at NDLAr}
The traditional deconvolution approach requires the full waveform and equally spaced samples. We developed algorithms to achieve these goals. We have a couple of charge records after the trigger starts and we compensate the missing waveform through a physics-guided template.

There are several important facts worth noting:
1. starts of waveform does not align with each other for waveforms with different amplitude at a fixed threshold. 2. the inferred full waveform is not the true full waveform. the assumption it is full waveform brings offset of starting point of charge. 3. the slow digitzation rate brigns further ambiguity. It is impossible to infer back when it is the best estimate of starting time of waveform below the digitzation frequency, while the trigger time is studied in a very fast manner.

The algorithm consists of three stages.

\subsubsection{Recovery of the full waveform}
The recovery of the full waveform is essentially performed in two passes. In the first pass, burst sequences are collected together as long as they are considered consecutive. a criterion is used, ADC-HOLD-DELAY, instead of ADC-CONVERSION-TIME.
Consecutive hits are merged into one sequences. It starts from the pre-triggering part, from the last hit to the trigger time, for charge at threshold, followed by recorded hits with threshold subtracted. If there are consecutive hits, it is appended to in sequence. We note that the time window for each hit are known except the first one.

In the second pass, the gap between sequences are filled by the template. We do a scan over te template. The sequence in the first pass are considered to record all ionization charge, then its relative height to the threshold provides a transit point on a single point charge. We scan over the template to idenfity the portion corresponding to that transition point in the time window between two consecutive sequences. For the first sequence, we direct use the the frist t0 to the transit point.

The time misalignment between sequences are treated in two separate way. 1. rounding them to the nearest global binning. 2. Enforce fractional exponent phase factor and transform sequences according to FFT.

In addition, the treatment of the detector dead time is also different. For the former one, charge lost in deadtime is asumed to be compensated by scaling the time window of the consecutive hits to one tick. For the second one, we do not compensate the deadtime.

\subsubsection{Decvonlution}
The recovered waveform are grouped into a big 3D array and deconvolved by the 3D field response, supplied with a Gaus filter.
\textcolor{red}{todo: deconvolution in blocks.}
\subsubsection{Identification of the active voxels}
The current strategy is to select voxels with charges above a threshold after a Gaus filter. The threshold is chosen to be 500e temporarily, standard deviation of the noise on frontend.
\textcolor{red}{todo: An ongoing analysis for identifying voxels using Wiener-inspired filter is performed.} The filter takes the form $F={\exp{-(\frac{\omega}{a}})^b}$ except $F(\omega=0)=0$.
\section{Performance evaluation in simulation}
\subsection{Simulation framework and setup}
The detector geometry is assumed to be the same as the ProtoDUNE-ND, also known as $2\times 2$ exerpeiment. A particle gun positron is taken as the input. We simulate the detector response using the modern GPU-accelerated package, TRED. 
\subsection{Performance evaluations}
\subsubsection{Choice of templates}
We have compared several different templates:
1. field respone function at the center of the collection pixel
2. average field response function at the collection pixel
3. average field response function for collection pixel plus its neighboring pixels (+1 pixels).
4. hybrid mode, average field response function for collection pixels on pixels with charge larger than 6000e-, average field response function for neighboring pixels with charge below than 6000e-.
5. \textcolor{red}{todo: An ongoing study using filtered template}.
\subsubsection{Choice of filters}
Several filters were tested. The filter along pixel axes is employed with a $\sigma_p~0.8$pitch. The filter along time axis is tested by several setups, starting from the one similar to the digitization rate, $1.6\mu s$, to several $\mu$s.
The comparisons

\subsubsection{Dependence on the threshold}
\textcolor{red}{todo: to compare the effects}
We have compared the setup with threshold-less, threshold at 1000e-, threshold at 5000e-. The results shows the slow digitzation rate 
\subsubsection{Dependence on the noise level}
We have compared the setups with and without noises simulated. The noise model is based on the LArPix-v2 chips. The noise is a subleading effects in our setup.

\section{Discussions}
\subsection{Truth bookkeeping from charge regression}
It is worth noting the charge measurements always suffers from detector smearing effects. Directly comparing the true ionization charge would inavoidaly suffers from the mismatch. Instead, we applies the same Gaus filter used in signal processing to the true ionization charge and compare them with the deconvolved charges. We define a set of time bins on active deconvolved charges and compare them with the true ionization charge. The amount of the true charge in the reconstructed area and time bin are then linked to the unique particle ID in this way. The simulation package tred saves the ionization charge along with the particle id. Pairing with te signal processing chain in this article, it provides the opportunitity for task which relies on the ionization history tracking, for instance, machine learning based reconstruction workflow.
\subsection{Hardware acceleration}
\textcolor{red}{todo: CPU vs. GPU performance benchmark}
The decvonlution through FFT is linearithmic and scales with the input well. The modern library, for instance, pytorch, usually benefits from the hardware acceleartion, for instance, on GPU.
\subsection{Requirements on front-end readout system}
It has been demonstrated that 1. short burst sequence 2. low threshold, are beneficial for the deconvolution as more information are preserved.
It is also worth demonstrating the uniformity of the threshold
\textcolor{red}{todo: various threshold and deconvolution performance?}
On the other hand, it is also worth noting, given the non-reset behaivor during the bursts, it is necessary to make sure the signal does not saturates.
\textcolor{red}{todo: maximum charge per burst sequence for some topologies}
\subsection{Shaping the induced signal by a shield grid (DUNE internal)}
According to the Ramo theorem, the adding a shield grid could essentially sharpen the signal in the collection pixel while suppress the long range induction on neighboring pixels. This improves the imaging sharply. The effective waveform become shorter so the missing waveform become smaller portion of the entire sequence. We demonstrate the study using the same setup, with and without a shield grid. The residue is shown. Much cleaner event display and smaller number of ghost hits are shown. The width of charge residue is also smaller.
\section{Summary}
We presented the signal processing at zero-suppressed detector. This method provides a way to unfold ionization charge from detector resolution effects, to quantify the detector resolution. It also provides a direct way to compare the true ionization charge with the reconstruced objects.

\appendix
\section{Internal documentations}
Here we put materials for internal DUNE review. We put spectra analyses here for reference. We need them to illustrate how the filter is chosen.





\acknowledgments

This is the most common positions for acknowledgments. A macro is
available to maintain the same layout and spelling of the heading.

\paragraph{Note added.} This is also a good position for notes added
after the paper has been written.


% Bibliography

%% [A] Recommended: using JHEP.bst file
%% \bibliographystyle{JHEP}
%% \bibliography{biblio.bib}

%% or
%% [B] Manual formatting (see below)
%% (i) We suggest to always provide author, title and journal data or doi:
%% in short all the informations that clearly identify a document.
%% (ii) please avoid comments such as "For a review'', "For some examples",
%% "and references therein" or move them in the text. In general, please leave only references in the bibliography and move all
%% accessory text in footnotes.
%% (iii) Also, please have only one work for each \bibitem.

\begin{thebibliography}{99}

\bibitem{a}
Author,
\emph{Title},
\emph{J. Abbrev.} {\bf vol} (year) pg.

\bibitem{b}
Author,
\emph{Title},
arxiv:1234.5678.

\bibitem{c}
Author,
\emph{Title},
Publisher (year).

\end{thebibliography}
\end{document}

